
\documentclass{article}
\usepackage{colortbl}
\usepackage{makecell}
\usepackage{multirow}
\usepackage{supertabular}

\begin{document}

\newcounter{utterance}

\centering \large Interaction Transcript for game `cladder', experiment `full\_v1.5\_default', episode 4050 with Qwen3.6{-}27B.
\vspace{24pt}

{ \footnotesize  \setcounter{utterance}{1}
\setlength{\tabcolsep}{0pt}
\begin{supertabular}{c@{$\;$}|p{.15\linewidth}@{}p{.15\linewidth}p{.15\linewidth}p{.15\linewidth}p{.15\linewidth}p{.15\linewidth}}
    \# & \multicolumn{2}{c}{Player} && \multicolumn{2}{c}{Game Master} \\
    \hline

    \theutterance \stepcounter{utterance}  
    & & & \multicolumn{4}{p{0.6\linewidth}}{
        \cellcolor[rgb]{0.9,0.9,0.9}{
            \makecell[{{p{\linewidth}}}]{
                \texttt{\tiny{[P1$\langle$GM]}}
                \texttt{You are an expert in causal inference. The following question is not a typical commonsense query, but rather a meticulously designed question created by a professor specializing in causal inference, intended to assess the students' mastery of the course content.} \\
\texttt{Imagine a self{-}contained, hypothetical world with only the following conditions, and without any unmentioned factors or causal relationships: Ability to speak english has a direct effect on diabetes and lifespan. Smoking has a direct effect on diabetes and lifespan. Diabetes has a direct effect on lifespan. Smoking is unobserved. For people who do not speak english, the probability of long lifespan is 77\%. For people who speak english, the probability of long lifespan is 45\%. Will speaking english decrease the chance of long lifespan?} \\
\texttt{Start your answer with 'Yes' or 'No', followed by additional reasoning or evidence to support your explanation.} \\
            }
        }
    }
     \\ \\

    \theutterance \stepcounter{utterance}  
    & \multicolumn{4}{p{0.6\linewidth}}{
        \cellcolor[rgb]{0.561,0.737,0.561}{
            \makecell[{{p{\linewidth}}}]{
                \texttt{\tiny{[P1$\rangle$GM]}}
                \texttt{No} \\
\\ 
\texttt{**Reasoning:**} \\
\\ 
\texttt{To determine whether speaking English *causes* a decrease in the chance of a long lifespan, we must distinguish between the observed association (correlation) and the causal effect. The problem provides observed probabilities:} \\
\texttt{{-} \$P(\textbackslash{}text\{Long Lifespan\} | \textbackslash{}text\{No English\}) = 0.77\$} \\
\texttt{{-} \$P(\textbackslash{}text\{Long Lifespan\} | \textbackslash{}text\{English\}) = 0.45\$} \\
\\ 
\texttt{At first glance, it appears that speaking English is associated with a lower probability of long lifespan. However, the problem explicitly states that **Smoking** is a variable that has a direct effect on both diabetes and lifespan, and crucially, **Smoking is unobserved**.} \\
\\ 
\texttt{In causal inference, if a variable (Smoking) affects the outcome (Lifespan) and is also associated with the treatment/exposure (Speaking English), it acts as a **confounder**. The problem description implies a structure where Smoking influences the outcome. While the prompt doesn't explicitly state that Smoking causes Speaking English, in typical causal inference problems of this type, the presence of an unobserved variable that affects the outcome and is potentially correlated with the exposure creates **confounding bias**.} \\
\\ 
\texttt{More specifically, let's look at the causal graph implied:} \\
\texttt{1.  Smoking \$\textbackslash{}rightarrow\$ Diabetes} \\
\texttt{2.  Smoking \$\textbackslash{}rightarrow\$ Lifespan} \\
\texttt{3.  Speaking English \$\textbackslash{}rightarrow\$ Diabetes} \\
            }
        }
    }
    & & \\ \\

    \theutterance \stepcounter{utterance}  
    & & & \multicolumn{2}{p{0.3\linewidth}}{
        \cellcolor[rgb]{0.9,0.9,0.9}{
            \makecell[{{p{\linewidth}}}]{
                \texttt{\tiny{[GM$|$GM]}}
                \texttt{No} \\
            }
        }
    }
    & & \\ \\

    \theutterance \stepcounter{utterance}  
    & & & \multicolumn{2}{p{0.3\linewidth}}{
        \cellcolor[rgb]{0.9,0.9,0.9}{
            \makecell[{{p{\linewidth}}}]{
                \texttt{\tiny{[GM$|$GM]}}
                \texttt{yes} \\
            }
        }
    }
    & & \\ \\

    \theutterance \stepcounter{utterance}  
    & & & \multicolumn{2}{p{0.3\linewidth}}{
        \cellcolor[rgb]{0.9,0.9,0.9}{
            \makecell[{{p{\linewidth}}}]{
                \texttt{\tiny{[GM$|$GM]}}
                \texttt{game\_result = LOSE} \\
            }
        }
    }
    & & \\ \\

\end{supertabular}
}

\end{document}
